\begin{ex}
Tính: $35 + 24$.
\shortans{59}
\loigiai{
Ta có $35 + 24 = 59$.
}
\end{ex}

\begin{ex}
Tính: $68 - 45$.
\shortans{23}
\loigiai{
Ta có $68 - 45 = 23$.
}
\end{ex}

\begin{ex}
Tính: $42 + 6$.
\shortans{48}
\loigiai{
Ta có $42 + 6 = 48$.
}
\end{ex}

\begin{ex}
Tính: $79 - 7$.
\shortans{72}
\loigiai{
Ta có $79 - 7 = 72$.
}
\end{ex}

\begin{ex}
Tính: $50 + 36$.
\shortans{86}
\loigiai{
Ta có $50 + 36 = 86$.
}
\end{ex}

\begin{ex}
Tính: $94 - 40$.
\shortans{54}
\loigiai{
Ta có $94 - 40 = 54$.
}
\end{ex}

\begin{ex}
Tính giá trị biểu thức: $32 + 15 - 24$.
\shortans{23}
\loigiai{
Ta có $32 + 15 = 47$, rồi $47 - 24 = 23$.
}
\end{ex}

\begin{ex}
Tính giá trị biểu thức: $87 - 35 + 21$.
\shortans{73}
\loigiai{
Ta có $87 - 35 = 52$, rồi $52 + 21 = 73$.
}
\end{ex}

\begin{ex}
Tính giá trị biểu thức: $20 + 30 + 40$.
\shortans{90}
\loigiai{
Ta có $20 + 30 = 50$, rồi $50 + 40 = 90$.
}
\end{ex}

\begin{ex}
Tính giá trị biểu thức: $90 - 40 - 20$.
\shortans{30}
\loigiai{
Ta có $90 - 40 = 50$, rồi $50 - 20 = 30$.
}
\end{ex}

\begin{ex}
Điền số thích hợp: $24 + \dots = 67$.
\shortans{43}
\loigiai{
Số cần điền: $67 - 24 = 43$.
}
\end{ex}

\begin{ex}
Điền số thích hợp: $89 - \dots = 34$.
\shortans{55}
\loigiai{
Số cần điền: $89 - 34 = 55$.
}
\end{ex}

\begin{ex}
Điền số thích hợp: $\dots + 31 = 75$.
\shortans{44}
\loigiai{
Số cần điền: $75 - 31 = 44$.
}
\end{ex}

\begin{ex}
Điền số thích hợp: $\dots - 23 = 52$.
\shortans{75}
\loigiai{
Số cần điền: $52 + 23 = 75$.
}
\end{ex}

\begin{ex}
Lớp 2A có 18 học sinh nam và 21 học sinh nữ. Hỏi lớp 2A có tất cả bao nhiêu học sinh?
\shortans{39}
\loigiai{
Cả lớp có: $18 + 21 = 39$ học sinh.
}
\end{ex}

\begin{ex}
Một cửa hàng nhập về 78 gói bánh, đã bán được 46 gói. Cửa hàng còn lại bao nhiêu gói bánh?
\shortans{32}
\loigiai{
Còn lại: $78 - 46 = 32$ gói bánh.
}
\end{ex}

\begin{ex}
Bác Ba thu hoạch được 54 bắp ngô nếp và 45 bắp ngô ngọt. Bác thu hoạch được tất cả bao nhiêu bắp ngô?
\shortans{99}
\loigiai{
Tổng số bắp: $54 + 45 = 99$ bắp ngô.
}
\end{ex}

\begin{ex}
Đàn gà có 68 con, bác nông dân bán đi 25 con. Đàn gà còn lại bao nhiêu con?
\shortans{43}
\loigiai{
Còn lại: $68 - 25 = 43$ con gà.
}
\end{ex}

\begin{ex}
Đoạn dây thứ nhất dài 45 cm, đoạn thứ hai dài 34 cm. Cả hai đoạn dây dài bao nhiêu xăng-ti-mét?
\shortans{79}
\loigiai{
Tổng chiều dài: $45 + 34 = 79$ cm.
}
\end{ex}

\begin{ex}
Sợi dây đồng dài 88 cm, cắt bớt đi 35 cm. Sợi dây đồng còn lại dài bao nhiêu xăng-ti-mét?
\shortans{53}
\loigiai{
Còn lại: $88 - 35 = 53$ cm.
}
\end{ex}

\begin{ex}
Số nào cộng với 32 thì được 76?
\shortans{44}
\loigiai{
Số đó là: $76 - 32 = 44$.
}
\end{ex}

\begin{ex}
Số nào trừ đi 25 thì được 41?
\shortans{66}
\loigiai{
Số đó là: $41 + 25 = 66$.
}
\end{ex}

\begin{ex}
Lấy 98 trừ đi số tròn chục lớn nhất có hai chữ số thì được kết quả là bao nhiêu?
\shortans{8}
\loigiai{
Số tròn chục lớn nhất có hai chữ số là 90. Hiệu: $98 - 90 = 8$.
}
\end{ex}

\begin{ex}
Tính tổng của số bé nhất có hai chữ số và số 75.
\shortans{85}
\loigiai{
Số bé nhất có hai chữ số là 10. Tổng: $10 + 75 = 85$.
}
\end{ex}

\begin{ex}
Điền số thích hợp vào chỗ chấm: $45 + 13 - \dots = 20$.
\shortans{38}
\loigiai{
Ta có $45 + 13 = 58$, vậy số cần điền là $58 - 20 = 38$.
}
\end{ex}
